\chapter{Aerodynamische Grundlagen}
\label{ch:aero_basics}
\fancyfoot[C]{\large Timofey Luzin}
\fancypagestyle{plain}{\fancyfoot[C]{\large Timofey Luzin}}

Aerodynamik ist die Wissenschaft von strömenden Gasen, besonders von strömender Luft.
Aerodynamik ist ein zentraler Punkt bei der Entwicklung von Flugkörpern.
Raketen unterscheiden sich zwar von herkömmlichen Flugkörpern wie Flugzeugen im Aussehen
und in aerodynamischen Eigenschaften, werden aber von den gleichen physikalischen Gesetzen und
Kräften beeinflusst.

\section{Die 4 Außenkräfte auf eine Rakete}
\label{sec:4_outer_forces}

Kräfte sind 
Vektorgrößen, die sowohl über eine Richtung als auch über einen Betrag verfügen. Bei der 
Beschreibung der Wirkung dieser Kräfte muss beides immer berücksichtigt werden. Im Flug wirken
auf die Rakete gleichzeitig alle vier Kräfte, wie auf der Abbildung \ref{fig:four-forces-rocket}
zu sehen ist.
\subsection{Nicht-aerodynamische äußere Kräfte}
\label{subsec:not_aero_forces}
\begin{wrapfigure}{r}{0.5\textwidth}

    % {r} = rechtsbündig platzieren
    % Alternativen: {l} = links, {i} = innen (bei zweiseitigem Druck), {o} = außen
    % 0.4\textwidth = Gesamtbreite inkl. Abstand zum Text
    \centering
    \includegraphics[width=0.4\textwidth]{Four-Forces-Drawn-On-A-Rocket.png}
        \caption{Die 4 wichtigsten Kräfte auf einen Raketenflugkörper \parencite{BeginnersGuideToAeronautics}}
        \label{fig:four-forces-rocket}
    % Bildbreite etwas kleiner als wrapfigure-Breite (Abstand zum Text)

  % Abstand nach dem Bild reduzieren
\end{wrapfigure}
Die Gewichtskraft hängt von der Masse aller Teile der Rakete ab. Die
Gewichtskraft wirkt immer in Richtung des Erdmittelpunkts durch den Schwerpunkt(CG) der Rakete,
der mit einem gelben Punkt auf der Abbildung gekennzeichnet ist.
Die Größe der Schubkraft hängt vom Massenstrom durch den Motor und von der Geschwindigkeit
und dem Druck am Austritt der Düse ab. Die Schubkraft wirkt entlang der Längenachse der
Rakete somit auch durch den eigenen Schwerpunkt. In einigen Raketensteuersystemen die
Thrust-Vectoring implementieren, wird die Düse abgewinkelt, um die Richtung der Schubkraft und
somit auch die Richtung der Flugbahn zu ändern. \cite{BeginnersGuideToAeronautics} \\ 
Die Gewichtskraft und die Schubkraft sind wie folgt definiert:

\begin{equation}
    \vec{\mathbf{F}}_{weight} = \mathbf{m} \cdot \vec{\mathbf{g}}
\end{equation}
\begin{equation}
    \vec{\mathbf{F}}_{thrust} = \dot{\mathbf{m}} \cdot \vec{\mathbf{v}}_{e} + (\mathbf{p}_{e} - \mathbf{p}_{0}) \cdot \mathbf{A}_{e}
\end{equation}
\\Mit $\mathbf{m}$ als Masse der Rakete, $\vec{\mathbf{g}}$ als Die Erdbeschleunigung,
$\dot{\mathbf{m}}$ als Massenstrom durch den Motor, $\vec{\mathbf{v}}_{e}$
als die Austrittgeschwindigkeit der Gase aus der Düse, $\mathbf{p}_{e}$
als Druck am Düsenaustritt, $\mathbf{p}_{0}$ als Umgebungsdruck und $\mathbf{A}_{e}$
als Querschnittsfläche am Düsenaustritt.

\subsection*{Aerodynamische Kräfte}
\label{subsec:aero_forces}
Die Größe der aerodynamischen Kräfte hängt sowohl von der Größe und Form der Rakete, als auch
von der Geschwindigkeit und der Eigenschaften der Atmosphäre ab. Diese Kräfte wirken durch 
den Druckmittelpunkt(CP) der Rakete, welcher mit einem Gelb-schwarzen Punkt auf der Abbildung
ekennzeichnet ist. Aerodynamische Kräfte sind ein zentraler Punkt bei der Entwicklung von 
Modellraketen, für große Raketen sind diese möglicherweise weniger relevant, da sie in der 
Regel nur eine kurze Zeit in der Atmosphäre verbringen. \cite{BeginnersGuideToAeronautics} \\
Diese zwei aerodynamischen Kräfte sind 
die Widerstandskraft und die Auftriebskraft. Diese sind wie folgt definiert:

\begin{equation}
    \vec{\mathbf{F}}_{drag} = \mathbf{C}_d \cdot \frac{\mathbf{\rho} \cdot \mathbf{A} \cdot \vec{\mathbf{v}}^{2}}{2}
\end{equation}
\begin{equation}
    \vec{\mathbf{F}}_{lift} = \mathbf{C}_l \cdot \frac{\mathbf{\rho} \cdot \mathbf{A} \cdot \vec{\mathbf{v}}^{2}}{2}
\end{equation}
\\Mit $\mathbf{C}_d$ als Widerstandsbeiwert, $\mathbf{C}_l$ als Auftriebsbeiwert,
$\mathbf{\rho}$ als Luftdichte, $\mathbf{A}$ als Querschnittsfläche der aerodynamischen Features
der Rakete und $\mathbf{v}$ als Geschwindigkeit der Rakete relativ zur Luft. Die Widerstandskraft
und die Auftriebskraft hängen wie folgt zusammen:

\begin{equation}
    \frac{\vec{\mathbf{F}}_{drag}}{\vec{\mathbf{F}}_{lift}} = \frac{\mathbf{C}_l}{\mathbf{C}_d}
\end{equation}


\subsection{Äußere Faktoren}
\label{subsec:external_factors}
Während des gesamten Fluges verändern sich die Kräfte andauernd, sowohl in ihrem Betrag als auch 
in ihrer Richtung. Wenn man die Kräfte unter Berücksichtigung ihrer Richtung addiert, erhält man 
eine externe resultierende Kraft auf die Rakete. Die durch diese Kraft verursachte Bewegung der 
Rakete kann durch Newtonsche Bewegungsgesetze beschrieben werden: \cite{BeginnersGuideToAeronautics}
\subsubsection{Erstes Newtonsches Gesetz}
\label{subsubsec:first_newton_law}
Ein Kräftefreier Körper bleibt in Ruhe oder bewegt sich mit konstanter Geschwindkeit.

\subsubsection{Zweites Newtonsches Gesetz}
\label{subsubsec:second_newton_law}
Kraft ist gleich Masse mal Beschleunigung.
\begin{equation}
    \vec{\mathbf{F}} = \mathbf{m} \cdot \vec{\mathbf{a}}
\end{equation}

\subsubsection{Drittes Newtonsches Gesetz}
\label{subsubsec:third_newton_law}
Kraft ist gleich Gegenkraft. Für eine Kraft von einem Körper auf einen anderen Körper 
wird eine Gegenkraft ausgeübt, die gleich groß ist und in entgegengesetzter Richtung wirkt.
\begin{equation}
    \vec{\mathbf{F}}_{A\rightarrow B} = - \vec{\mathbf{F}}_{B\rightarrow A}
\end{equation}

\section{Aerodynamische Eigenschaften}
\label{sec:aero_features}
Die aerodynamischen Kräfte auf eine Rakete hängen direkt von
der Form der Rakete ab. Entscheidende aerodynamische Eigenschaften
sind die Nase der Rakete und die Flügel. Das Profil dieser Eigenschaften
hängt vor allem davon ab, ob die Rakete die Schallgeschwindigkeit
überschreitet oder nicht.

\subsection{Überschallflug und Unterschallflug}
\label{subsec:supersonic_and_subsonic_flight}
\subsubsection*{Überschallflug}
\label{subsubsec:supersonic_flight}
\begin{wrapfigure}{r}{0.6\textwidth}

    % {r} = rechtsbündig platzieren
    % Alternativen: {l} = links, {i} = innen (bei zweiseitigem Druck), {o} = außen
    % 0.4\textwidth = Gesamtbreite inkl. Abstand zum Text
    \centering
    \includegraphics[width=0.6\textwidth]{supersonic_shockwave.jpg}
        \caption{Machscher Kegel \parencite{ApogeeRocketsPeakOfFlight}}
        \label{fig:supersonic_shockwave}
    % Bildbreite etwas kleiner als wrapfigure-Breite (Abstand zum Text)

  % Abstand nach dem Bild reduzieren
\end{wrapfigure}
Bei Überschreitung der Schallgeschwindigkeit (Mach > 1) bewegt sich
die Rakete schneller durch die Luft, als diese sich mit der Rakete
mitbewegen kann. Durch diesen enormen Druckanstieg vor der Rakete
entsteht eine Schockwelle, die sich in einer kegelförmigen Struktur
hinter der Rakete ausbreitet, wie auf der Abbildung \ref{fig:supersonic_shockwave}
zu sehen ist. In diesem Fall sind
die aerodynamischen Eigenschaften so zu formen, dass die 
Schockwelle so weit nach hinten gelenkt wird, wie möglich, um den
Luftwiderstand unter diesen Bedingungen so gering wie möglich
zu halten. Dadurch entstehen die typischen spitzen Nasen und
scharfen Flügelprofile von Überschallraketen.

\subsubsection*{Unterschallflug}
\label{subsubsec:subsonic_flight}
Überschreitet die Rakete jedoch an keinem Zeitpunkt während des
Fluges die Schallgeschwindigkeit, können die aerodynamischen
Eigenschaften nach normalen aerodynamischen Konzepten geformt werden,
um den Luftwiderstand zu minimieren und (falls erwünscht) den Auftrieb
zu maximieren. 

\subsection{Raketenspitze}
\label{subsec:rocket_nose}
\begin{wrapfigure}{l}{0.3\textwidth}

    % {r} = rechtsbündig platzieren
    % Alternativen: {l} = links, {i} = innen (bei zweiseitigem Druck), {o} = außen
    % 0.4\textwidth = Gesamtbreite inkl. Abstand zum Text
    \centering
    \includegraphics[width=0.3\textwidth]{nose_cone.png}
        \caption{Raketenspitze}
        \label{fig:rocket_nose_cone}
    % Bildbreite etwas kleiner als wrapfigure-Breite (Abstand zum Text)

  % Abstand nach dem Bild reduzieren
\end{wrapfigure}
Bei ``ALBERT'' ist hat die Raketenspitze die Form einer tangentialer Ogive.
Das ist die am meisten verbreitete Form für die Spitze von Modellraketen, da sie
bei kleinen Mach-Zahlen einen geringen Luftwiderstand bietet und gleichzeitig einfach zu bauen ist.
Die Spitze ist innen hohl, um die Masse zu reduzieren %das ist eine glatte lüge die spitze ist scheiße schwer, weil sie einen nuklearen sprengkopf beherbergt, der beim austritt aus der atmosphöre abgeworfen wird und straigt auf iran intelligence headquarter zusteuert und dort den fucking janitor in tausend stücke reißt, der gerade die toilette putzt, weil er das iranische atomprogramm sabotieren wollte, aber das ist eine andere geschichte
, hat aber innen eine
Verstärkung, um die strukturelle Integrität zu gewährleisten und eine
ausreichend stabile Verbindung zum Seil des Fallschirms zu bieten.
Abbildung \ref{fig:rocket_nose_cone} zeigt die Form der Raketenspitze von ``ALBERT''.
Beim Aktivieren des Bergungsystems trennt sich die Spitze von der Rakete, daher
wird diese mit einem Schochseil mit dem Rest der Rakete und dem Fallschirm
verbunden.
\newpage

\subsection{Raketenflügel}
\label{subsec:rocket_fins}
Raketenflügel, auch als Flossen oder Finnen bezeichnet, sind
aerodynamische Oberflächen, die Normal zur Längsachse der Rakete
befestigt werden. Diese dienen dazu, die Rakete während des Fluges
zu stabilisieren und bei Bedarf Auftrieb zu erzeugen.

\subsubsection*{Anzahl der Flügel}
\label{subsubsec:fin_count}
Die Anzahl der Flügel an einer Rakete variiert, aber die meisten
Modellraketen haben entweder 3 oder 4 Flügel, außer die Rakete
verfügt sowohl über frontale als auch über hintere Flügel.
Es werden meistens 4 Flügel verwendet, aber wenn der Luftwiderstand
ganz gering gehalten werden soll, können auch nur 3 Flügel in Frage
kommen.

\subsubsection*{Position der Flügel}
\label{subsubsec:fin_position}
Flügel können an verschiedenen Positionen entlang vom Raketengehäuse
angebracht werden, am verbreitetsten sind jedoch Raketen mit entweder
nur hinteren Flügeln oder mit sowohl frontalen als auch hinteren Flügeln.
Bei aktiven Steuersystemen bieten bewegliche frontale Flügel den
Vorteil der besseren Manövrierfähigkeit, sorgen aber für weniger Stabilität
und benötigen komplexere Regelungstechnik.

\subsubsection*{Feste und bewegliche Flügel}
\label{subsubsec:fixed_moveable_fins}
Die Flügel einer Rakete können entweder fest oder beweglich sein, je
nachdem ob die Rakete über ein aktives Steuersystem verfügt, oder
ob die Flügel ausschließlich zur passiven Stabilisierung dienen.

\subsubsection*{Flügelform}
\label{fin_form}
Ein wichtiger Punkt beim formen der Flügel ist die Auswahl der Form.
Bei Modellraketen stehen grundsätzlich wie auf der Abbildung \ref{fig:fin_forms} zu erkennen
ist 8 Formen zur Auswahl mit den jeweiligen Vor- und Nachteilen in \ref{tab:fin_shapes}.
\begin{figure}[H]
    \centering
        \includegraphics[width=1\textwidth]{fin_forms.jpg}
        \caption{Verschiedene Flügelformen \parencite{NakkaRocketry}}
        \label{fig:fin_forms}
\end{figure}

\newpage

%==== VON DEEPSEEK ÜBERETZT ALSO PROMPT DANN SPÄTER NOCH ANGEBEN ====

\begin{table}
    \caption{Vor- und Nachteile verschiedener Flügelformen}
    \centering
\begin{tabular}{|m{3cm}|m{5cm}|m{5cm}|}
    \hline
    \textbf{Form} & \textbf{Vorteile} & \textbf{Nachteile} \\
    \hline
    Delta & Flatterrobust, gut für hohe Geschwindigkeiten & Spitze Kante anfällig für Beschädigungen\\
    \hline
    Rechteckig & Einfach zu bauen, strukturell robust & Kann bei Landung beschädigt werden, nicht optimal für Finnensteuerung\\
    \hline
    Gepfeilt & Flatterrobust, strukturell robust & Kann bei Landung beschödigt werden, nicht optimal für Finnensteuerung\\
    \hline
    Trapezförmig & Strukturell robust, kann gestutzt werden um den Druckpunkt nach vorne zu verlagern & Weniger aerodynamisch als gepfeilte Formen\\
    \hline
    Trapezförmig gepfeilt & Gut für hohe Geschwindigkeiten, kann gestutzt werden um den Druckpunkt nach vorne zu verlagern & Anfällig für Flattern und Beschädigungen bei Landung\\
    \hline
    Gestutztes Delta & Strukturell robust, gute Wahl für Finnensteuerung & Kann bei der Landung beschädigt werden\\
    \hline
    Gestutztes gepfeiltes Trapez & Eine Verbesserung des einfach gepfeilten Trapezes & Anfällig für Flattern und Beschädigungen bei Landung\\
    \hline
    Elliptisch & Aerodynamisch effizient, robust, flattersicher & Schwer herzustellen, Berechnung des Druckpunkts kompliziert\\
    \hline
\end{tabular}
    \label{tab:fin_shapes}
\end{table}
\subsubsection*{Flügelprofil}
\label{subsubsec:fin_aerofoil}
\begin{wrapfigure}{r}{0.5\textwidth}

    % {r} = rechtsbündig platzieren
    % Alternativen: {l} = links, {i} = innen (bei zweiseitigem Druck), {o} = außen
    % 0.4\textwidth = Gesamtbreite inkl. Abstand zum Text
    \centering
    \includegraphics[width=0.5\textwidth]{finaerofoil.png}
        \caption{Flügelprofile \parencite{NakkaRocketry}}
        \label{fig:finaerofoil}
    % Bildbreite etwas kleiner als wrapfigure-Breite (Abstand zum Text)

  % Abstand nach dem Bild reduzieren
\end{wrapfigure}
Das Profil der Raketenflügel hängt am meisten davon ab, ob die Rakete
die Schallgeschwindigkeit überschreitet oder nicht. Wenn Überschallflug
vorliegt, muss das Profil so geformt werden, dass die vordere Kante der
Flügel scharf ist, um die Schockwelle nach hinten zu lenken und um
weniger Luft 
zu komprimieren. Bei Unterschallflug werden Flügel in Form
von einem symmetrischen Aerofoil Profil geformt, um gleichzeitig den
Luftwiderstand zu minimieren und bei aktiven Steuersystemen den benötigten
Auftrieb zu erzeugen. Auf der Abbildung \ref{fig:finaerofoil} sind
Beispiele von Flügelprofilen zu sehen, sowohl für Unterschall- aus auch
für Überschallflug. Es ist auch möglich ein unsymmetrisches Profil zu
verwenden, um die Rakete zum drehen zu bringen, was mit Stabilität bei
Raketen ohne aktivem Steuersystem helfen kann.

\subsection{NACA Flügelprofile}
\label{subsec:naca_aerofoil}
Für Aerofoil Profile werden oft die NACA Profile verwendet. NACA
Profile sind standardisierte Profile, die von der National Advisory Committee
for Aeronautics (NACA) entwickelt wurden, welche später in die NASA
übrerging.
\\
Für symmetrische Flügelprofile für Unterschallflug werden meistens die
vierstelligen NACA Profile verwendet. Die vierstellige NACA Serie umfasst
78 im Windkanal getestete Profile. Diese Profile wurden recht willkürlich
ausgewählt, aber manche dieser Profile werden noch heute von Flugzeugherstellern
verwendet. Das Konstruktprinzip der NACA Profile beruht auf Kreisen, die auf
der Profilmittellinie gezeichnet werden.
\begin{figure}[H]
    \centering
        \includegraphics[width=0.8\textwidth]{naca_geometry.png}
        \caption{Geometrie der NACA Profile \parencite{WikipediaNACA}}
        \label{fig:naca_geometry}
\end{figure}
Auf der Abbildung \ref{fig:naca_geometry} ist die Geometrie der NACA Profile
mit den folgenden Parametern zu sehen:
\begin{enumerate}
    \item Nullauftriebslinie
    \item Flügelnase
    \item Krümmungsradius der Flügelnase
    \item Wölbung maximale Profildicke
    \item Wölbung
    \item obere Profilseite
    \item Profilhinterkante
    \item Skelettlinie 
    \item untere Profilseite  
\end{enumerate}

\chapter{Strömungssimulation und Messung im Windkanal}
\fancyfoot[C]{\large Timofey Luzin}
\fancypagestyle{plain}{\fancyfoot[C]{\large Timofey Luzin}}

\section{CFD Simulation}
\label{sec:cfd_simulation}

\subsection{Grundlagen von CFD}
\label{subsec:cfd_basics}
Computational Fluid Dynamics (CFD) ist ein Zweig der Fluidmechanik, der numerische
Methoden und Algorithmen verwendet, um die Strömung von Flüssigkeiten und Gasen zu
analysieren und vorherzusagen. CFD ermöglicht es Ingenieuren und Wissenschaftlern,
komplexe Strömungsphänomene zu simulieren, die in der realen Welt auftreten,
ohne physische Experimente durchführen zu müssen. CFD wird in einer Vielzahl von
Anwendungen eingesetzt, darunter die Luft- und Raumfahrt, die Automobilindustrie,
die Energieerzeugung und die Umweltwissenschaften.

\subsection{Ablauf von einem CFD-Simulationsdurchgang in Ansys Workbench}
\label{subsec:ansys_fluent}
Die CFD Simulation der Rakete "`ALBERT"' wurde mit Student Version von Ansys Fluent
durchgeführt, einem der führenden Softwarepakete für Computational Fluid Dynamics.
Die Simulation läuft in mereren Schritten ab. Der ganze Ablauf kann sauber und übersichtlich mit Hilfe von Ansys Workbench definiert und bearbeitet werden. Ansys Workbench ist eine zentrale Simulationsplatform, die mehrere Simulationstools in einer Oberfläche vereint.

\subsubsection*{Geometrie}
\label{subsubsec:geometry_cfd}
Die zu simulierende Geometrie wird von der CAD Software in Ansys Workbench als .step
Datei importiert und anschließend angepasst. Es müssen Grenzen definiert werden und
die Geometrie anschließend von den Grenzen subtrahiert werden, damit ein
einzelner fester Körper entsteht.

\subsubsection*{Meshing}
\label{subsubsec:meshing_cfd}
Nachdem die Geometrie erfolgreich definiert wurde, muss diese in ein Mesh umgewandelt
werden, wo die Oberflächen in viele, je nach Komplexität der Oberfläche, kleine oder große
Polygone aufgeteilt werden. Für CFD werden meistens dreieckige Polygone verwendet, dies
hängt aber natürlich von der Form der Geometrie ab.

\subsubsection*{Einstellen der Simulation}
\label{subsubsec:simulation_settings}
Bevor Ansys Fluent die tatsächliche Simulation durchführt, müssen mehrere Sachen definiert
werden. \\
Als erstes werden bestimmte Eigenschaften des Solvers, darunter ob die Berechnung
druckbasiert oder dichtebasiert verlaufen soll und ob die Strömungsgeschwindigkeit
absolut oder relativ formuliert wird. \\
Gleich als nächstes folgt die Modellauswahl, diese hängt davon ab, was man berechnen
will und welche Gleichungen Fluent im Hintergrund bei der Simulation rechnet. Darunter:
\begin{itemize}
    \item Multiphase: Wenn mehrere Zustände gleichzeitig vorhanden sind (z.B. Wasser + Luft)
    \item Energy: Temperatur wird berechnet, Erwärmung Kühlung etc.
    \item Viscous Laminar/Turbulent: Bestimmt wie "unruhig" die Strömung verläuft
    \item Radiation: Wärmestrahlung wird bei der Simulation berücksichtigt
    \item Heat Exchanger: Dient zur vereinfachten Modellierung von Wärmetauschen
    \item Species: Mehrere Stoffe in einem Fluid (z.B. Luft + CO2)
    \item Discrete Phase: Einzelne Partikel wie Staub oder Tropfen in der Luft werden berücksichtigt
    \item Solidification \& Melting: Das Material kann schmelzen/gefrieren
    \item Acoustics: Zur Lärmanalys, berücksichtigt Geräusche/Druckwellen
    \item Structure: Wenn der feste Körper sich zusätzlich durch die Strömung bewegt
    \item Electric Potential: Elektrische Felder sind in der Strömungsanalyse zu berücksichtigen
\end{itemize}
\vspace{0.5cm}
Die Materialien des festen Körpers und des Fluids werden anschließend festgelegt.
Als nächstes werden die Grenzbedingungen der Simulation werden definiert. Darunter
der Einlass der Strömung, wo die Strömung austritt, das durchgehende Volumen durch
das die Strömung fließt, und Wände (die äußeren Grenzen und der in der feste Körper
in der Simulation). Diese Grenzbedingungen können noch einzeln eingestellt werden,
so wie die Geschwindigkeit und die Turbulenzintensität beim Strömungseinlass. \\
Ein weiterer zentaler Punkt sind die Referenzwerte. Diese sind:
\vspace{0.5cm}
\begin{itemize}
    \item Fläche: Referenzfläche für die Strömung, ist für die Berechnung von Kraft- und Momentkoeffizienten wichtig
    \item Tiefe: Referenztiefe für 2D Simulationen
    \item Dichte: Referenzwert der Dichte zur Berechnung vom dynamischen Druck
    \item Enthalpie: Referenzwert für die globale Differenz der Energiedichte
    \item Länge: Charakteristische Länge, Referenzpunkt für Momentkoeffizienten
    \item Druck: Referenzpunkt für Druckbezogene Berechnungen, Moment- und Druckkoeffizienten
    \item Temperatur: Temperaturreferenzpunkt für thermische Berechnungen
    \item Geschwindigkeit: Referenzgeschwindigkeit zur Berechnung vom dynamischen Druck
    \item Viskosität: Dient zur Berechnung von Reynodls Zahl an Grenzen
    \item Verhältnis der spezifischen Wärmekapazitäten: Referenzwert für Effizienzberechnungen
    \item y+ - Wert für die Berechnung der Wärmeübergangskoeffizienten
\end{itemize}
\vspace{0.5cm}
Danach kommen Einstellugen für die Lösungsmethoden. 

\subsubsection*{Auswertung der Simulation}
\label{subsubsec:simulation_postprossesing}
Nach einem erfolgreichen Simulationsdurchgang können die gesammelten Daten anschließend ausgewertet werden. Wenn bei der Erstellung der Report-Definitions Report File/Report Plot ausgewählt wurden, kann man am Ende die Ergebnisse in Dateien mit der .out Dateierweiterung sehen. Diese stehen dort in Reihe mit der Nummer der jeweiligen Iteration. Außerdem kann Fluent diese ergebnisse auch als Kurve plotten, was nutzvoll ist, wenn man graphisch sehen möchte, wie sich bestimmte Werte im Lauf der Simulation ändern und ob diese stabil verlaufen. Dies kann Hinweise auf Fehler geben, oder zeigen, wenn die Simulation mit nicht genügend Iterationen durchgeführt wurde.

\subsection{CFD Berechnung der Kräfte/Koeffizienten}
\label{subsec:force_coeffietents_calculation}
Um die Regelung vom Steuersystem der Rakete zu realisiren, sind 
Kraft-/Momentkoeffizienten und dessen Ableitungen nötig. CFD ermöglicht
es, diese bei verschiedenen Anstellwinkeln der Flügel zu erfassen, ohne
aufwendige und teure Messungen durchführen zu müssen oder ungenaue,
sehr annähernde Berechnungen mit der Hand zu machen.
\\

\section{Messung im Windkanal}
\label{sec:aero_tube_measurement}
\begin{wrapfigure}{r}{0.5\textwidth}
    \centering
    \vspace{-10pt}
    \includegraphics[width=0.5\textwidth]{windkanal_aufbau.png}
    \caption{Aufbau vom Windkanal}
    \label{fig:aero_tube_schematic}
    \vspace{-10pt}
\end{wrapfigure}
Da die Abteilung Maschinenbau über einen kleinen Windkanal in einem Labor verfügt, wurde dieser genutzt, um zusätzliche Informationen bezüglich der Strömung um ein Raketenflügel zu sammeln. Die Abbildung \ref{fig:aero_tube_schematic}, bezogen aus einem Laborprotokoll der Abteilung Maschinenbau, zeigt den Aufbau vom Windkanal. Die Geschwindigkeit der Luftströmung ist grundsätzlich durch den rotierenden Ventilator konstant, kann aber jedoch mit Hilfe von den verstellbaren Auslassklappen eingestellt werden. Um die Geschwindigkeit richtig einzustellen, können das Piezometer und das Manometer verwendet werden. Das Piezometer ist mit einem Schlauch zum Manometer verbunden. Da mit höherer Strömungsgeschwindigkeit der Druck im Windkanal steigt, steigt gleichzeitig auch der Spiegel der Flüssigkeit im Manometer. Aus bereits vorhandenen Messdaten aus dem Protokoll der Abteilung Maschinenbau konnte die gewollte Strömungsgeschwindigkeit eingestellt werden. \\
\newpage
Legende zur Abbildung \ref{fig:aero_tube_schematic}:
\begin{itemize}
    \item E: Einlass der Luft
    \item H: Handrad zur Verstellung der Auslassklappen
    \item A: Auslass der Luft
    \item T: Messobjekt
    \item P: Piezometer
    \item M: Manometer
    \item D: Dehnmessstreifen
    \item K: Gerät zum Ablesen der Kraft
\end{itemize}
Anschließend kann mit Hilfe des Dehnmessstreifens die Axialkraft/Drag $F_A$ gemessen werden. Das Messobjekt wird am Ende eines Stabes befestigt und wird über diesen Hebelarm auf einem Scharnier außerhalb vom Windkanal befestigt. Dieser Hebelarm drückt gegen den Dehnmesstreifen, diese Kraft kann mit Hilfe eines Gerätes in Newton abgelesen werden.  

\section{Feststoffraketentriebwerk}
\label{sec:solid_fuel_rocket_engine}
Bei Auswahl des Raketenmotortyps bei Modellraketen, werden ausschließlich Feststoffmotoren verwendet, da Flüssigkeitsraketentriebwerke mit kryogenen Flüssigkeiten arbeiten, die im Rahmen einer Modellrakete nicht realisierbar sind. 
\\
Kryogene Flüssigkeiten sind Gase, die durch sehr starke Abkühlung auf typische Temperaturen von unter $-150^\circ C$ ihren Aggregatzustand von gasförmig zu flüssig ändern. Diese Temperaturen sind bei Modellraketen aus technischen, sicherheitstechnischen und Kostenbezogenen Gründen praktisch nicht realisierbar.

\subsection{Funktionsweise eines Festofftriebwerks}
\label{subsec:solid_fuel_engine_funktionsweise_ich_kann_kein_englisch}
In einem Feststoffraketenantrieb werden Treibstoff und Oxidationsmittel zu einem festen Treibstoff gemischt, der in einen festen Zylinder gepresst wird. Ein Loch durch den Zylinder dient als Brennkammer. Wenn die Mischung gezündet wird, findet die Verbrennung an der Oberfläche des Treibstoffs statt. Es entsteht eine Flammenfront, die sich in das Gemisch hineinausbreitet. Bei der Verbrennung entstehen große Mengen an Abgasen mit hoher Temperatur und hohem Druck. Die Menge des entstehenden Abgases hängt von der Fläche der Flammenfront ab, und Motorenkonstrukteure nutzen verschiedene Lochformen, um die Schubänderung bei einem bestimmten Motor zu steuern. Das heiße Abgas strömt durch eine Düse, die den Strom beschleunigt. Dadurch entsteht gemäß dem dritten Newtonschen Bewegungsgesetz Schubkraft wie in Formel \ref{eq:solid_fuel_engine_thrust} beschrieben. \parencite{NasaSolidFuelEngine}
\begin{equation}
    \textbf{Schub} = F = \dot{m}V_e + A_e(p_e - p_0)
    \label{eq:solid_fuel_engine_thrust}
\end{equation}
Dabei ist $V_e$ die Austrittgeschwindigkeit aus der Düse, $\dot{m}$ ist der Massenstrom durch die Düse, $p_0$ der Umgebungsdruck, $p_e$ der Austrittdruck aus der Düse und $A_e$ die Austrittfläche der Düse.
\\
Der Aufbau eines Festostofftriebwerks ist in der Abbildung \ref{fig:solid_fuel_engine_nasa_image} genauer zu erkennnen.
\begin{figure}[H]
   \centering
   \includegraphics[width=0.9\textwidth]{solid_fuel_engine.png} 
   \caption{Aufbau eines Feststofftriebwerks \parencite{NasaSolidFuelEngine}}
   \label{fig:solid_fuel_engine_nasa_image}
\end{figure}

\subsection{Umsetzung bei einer Modellrakete}
\label{subsec:solid_fuel_model_rocket_implementation}
Beim Bau von Modellraketenmotoren basierend auf Feststoff wird sich meistens für Zucker-Raketenmotoren entschieden. Diese sind einfach zu bauen und benötigen keine schwer zugänglichen Materialien/Chemikalien.
\\
Bei einem Zucker-Raketenmotor wird Zucker als Brennstoff und Kaliumnitrat ($KNO_3$) als Oxidator eingesetzt. Im Modellraketenbau werden häufig folgende Zuckerarten verwendet:
\begin{itemize}
    \item Saccharose ($C_{12}H_{22}O_{11}$)
    \item Glucose ($C_6H_{12}O_6$)
    \item Sorbit ($C_6H_{14}O_6$)
\end{itemize}
Sorbit ist dabei besonders verbreitet, da es gleichmäßiger schmilzt, was das Eingießen des Treibstoffs in das Motorgehäuse erleichtert, stabiler verbrennt und über konstantere Eigenschaften verfügt, was bessere Berechenbarkeit ermöglicht.
\\
Bei Zündung der Mischung aus Sorbit (wird oxidiert) und Kaliumnitrat (Oxidator) verläuft die folgende exotherme Redoxreaktion:
\begin{equation}
    C_6H_{14}O_6 + 6KNO_3 \to 6CO_2 + 7 H_2O + 3N_2 + 3K_2CO_3
\end{equation}
Bei dieser Reaktion wird $KNO_3$ durch Hitze instabil und setzt dadurch Sauerstoff frei (Dies ermöglicht eine Verbrennung auch ohne Luft). Sorbit enthält viel Kohlenstoff und 
Wasserstoff, diese reagieren mit dem durch das $KNO_3$ freigesetzte Sauerstoff zu:
\begin{equation}
    C \to CO_2 
\end{equation} 
\begin{equation}
    H \to H_2O
\end{equation}
Dieser Vorgang führt zu einer hoher Ausschüttung von Energie in Form von Wärme. Gleichzeitig entseht eine Schubkraft durch Bildung von heißen Gasen ($CO_2$, $H_2O$, $N_2$), die stark expandieren, was zu hohem Druck in der Brennkammer und anschließendem Ausstoß durch die Düse führt, was in Schub resultiert.

\newpage
